\chapter{Cystoscopy}

\section*{Definition and Overview}
Cystoscopy is a procedure to look inside the bladder and urethra using a thin tube with a camera called a cystoscope. It allows a urologist to examine the lining of the bladder, diagnose conditions such as bladder stones, tumours, and inflammation, and perform biopsies or small treatments. Cystoscopy is usually done as a day procedure under local anaesthetic or sedation, and it takes only a few minutes. It is a safe and common investigation, and most people are able to go home the same day.

\section*{Epidemiology}
Cystoscopy is one of the most common procedures performed by urologists. It is used to investigate blood in the urine, frequent urinary tract infections, bladder pain, and urinary symptoms that do not respond to treatment, and to follow up people who have had bladder tumours. Bladder cancer, which cystoscopy detects, is one of the more common cancers, particularly in older men who smoke. Flexible cystoscopy under local anaesthetic is the standard approach for diagnosis, with rigid cystoscopy used for procedures under anaesthetic.

\section*{Aetiology and Pathophysiology}
Cystoscopy is performed to diagnose and treat problems within the bladder and urethra.

\begin{itemize}
    \item \textbf{Haematuria:} blood in the urine, which requires examination of the bladder lining
    \item \textbf{Bladder tumours:} the main reason for surveillance cystoscopy
    \item \textbf{Bladder stones:} hard deposits that can cause pain and bleeding
    \item \textbf{Recurrent urinary tract infections:} to look for underlying causes
    \item \textbf{Bladder pain and interstitial cystitis:} to assess the bladder lining
    \item \textbf{Urinary obstruction:} to inspect the prostate and urethra
    \item \textbf{The procedure:} the cystoscope is passed through the urethra into the bladder, and water is used to distend the bladder for a clear view
\end{itemize}

\section*{Clinical Features}
Cystoscopy is a test, so the features are the symptoms that prompt it.

\begin{itemize}
    \item Blood in the urine, visible or microscopic
    \item Frequent, urgent, or painful urination
    \item Recurrent urinary tract infections
    \item Bladder pain or discomfort
    \item Difficulty passing urine or a weak stream
    \item A previous diagnosis of bladder tumour requiring follow-up
    \item Abnormal urine tests or imaging findings
\end{itemize}

\section*{Differential Diagnosis}
The procedure answers specific questions and is combined with other tests.

\begin{itemize}
    \item Urinary tract infection, diagnosed by urine tests
    \item Bladder cancer, detected or excluded by cystoscopy
    \item Bladder stones and other structural problems
    \item Prostate disease in men, assessed with ultrasound and examination
    \item Interstitial cystitis, a diagnosis made when other causes are excluded
    \item Neurological or functional causes of urinary symptoms
\end{itemize}

\section*{When to Seek Medical Care}
A urologist refers people for cystoscopy when symptoms warrant it. After the procedure, contact the team if you develop significant problems.

\textbf{Call 000 (or local emergency services) for:}
\begin{itemize}
    \item The inability to pass urine after the procedure
    \item Heavy bleeding or passing clots
    \item Fever, chills, or severe abdominal pain (possible infection)
    \item Severe pain that is not settling
\end{itemize}

\section*{Diagnosis and Investigations}
Cystoscopy is often combined with other investigations.

\begin{itemize}
    \item \textbf{Urine tests:} dipstick, microscopy, culture, and cytology
    \item \textbf{Ultrasound:} of the kidneys and bladder, often before cystoscopy
    \item \textbf{CT urogram:} detailed imaging of the urinary tract
    \item \textbf{Flexible cystoscopy:} under local anaesthetic, the standard diagnostic test
    \item \textbf{Rigid cystoscopy:} under general or spinal anaesthetic, when biopsy or treatment is planned
    \item \textbf{Biopsy:} tissue samples taken during the procedure for analysis
    \item \textbf{Follow-up:} regular surveillance cystoscopy after bladder tumour treatment
\end{itemize}

\section*{Management}
Management depends on what the procedure finds.

\begin{itemize}
    \item \textbf{Normal findings:} reassurance, with treatment of any underlying condition
    \item \textbf{Bladder tumour:} removal (resection) through the cystoscope, with further treatment based on the biopsy result
    \item \textbf{Bladder stones:} crushing and removal during cystoscopy
    \item \textbf{Infection:} antibiotics
    \item \textbf{Interstitial cystitis:} specialist management, including bladder instillation therapies
    \item \textbf{Strictures:} dilatation or incision of a narrowed urethra
    \item \textbf{Follow-up:} surveillance schedules for people with previous tumours
\end{itemize}

\section*{Complications}
Cystoscopy is safe, but complications can occur.

\begin{itemize}
    \item Temporary burning or blood in the urine, common in the first day or two
    \item Urinary tract infection
    \item Difficulty passing urine after the procedure
    \item Injury to the urethra or bladder (rare)
    \item Bleeding, usually minor
    \item Pain or discomfort
    \item Anaesthetic-related risks for procedures under general anaesthetic
\end{itemize}

\section*{Prognosis and Outcomes}
Cystoscopy is a quick, safe procedure with rapid recovery, and most people return to normal activities within a day or two. It provides a definitive view of the bladder, allowing accurate diagnosis and treatment. The outcomes depend on what is found: normal findings are reassuring, bladder stones are usually cured by removal, and bladder tumours are managed according to their stage and grade, with regular surveillance. Mild post-procedure burning and blood in the urine settle within a few days.

\section*{Prevention and Self-Care}
\begin{itemize}
    \item Follow the preparation instructions, including any fasting requirements
    \item Tell the urologist about medicines, especially blood thinners, allergies, and pregnancy
    \item Drink extra fluids after the procedure
    \item Expect mild burning and blood in the urine for a day or two
    \item Take pain relief as advised
    \item Contact the team for fever, heavy bleeding, or inability to pass urine
    \item Attend follow-up appointments, including surveillance cystoscopy
\end{itemize}

\section*{Sources and Further Reading}
The following reputable sources provide further information for healthcare students and patients seeking reliable, current advice about cystoscopy.

\begin{itemize}
    \item Healthdirect Australia. \textit{Cystoscopy.} https://www.healthdirect.gov.au/
    \item Urological Society of Australia and New Zealand. \textit{Cystoscopy patient information.} https://www.usanz.org.au/
    \item Inside Radiology (Royal Australian and New Zealand College of Radiologists). \textit{Cystoscopy and imaging of the urinary tract.} https://www.insideradiology.com.au/
    \item NHS. \textit{Cystoscopy.} https://www.nhs.uk/conditions/cystoscopy/
    \item Mayo Clinic. \textit{Cystoscopy: overview and what to expect.} https://www.mayoclinic.org/tests-procedures/cystoscopy/
    \item MedlinePlus. \textit{Cystoscopy.} https://medlineplus.gov/ency/article/003903.htm
\end{itemize}
